%% ============================================================
%% 片段：s03c-topology   对应原文 PDF 第 62--73 页
%% 本文件由 gw/convert.py 从 arXiv 源码（hep-th/0612073v2）自动生成：
%%   * 公式/编号/引用/图 已经锁定为原文内容，请勿改动任何数学内容与编号；
%%   * 只需把英文正文译为中文，并在每节开头加“本节重点”、关键处加译注。
%% ============================================================


\subsection{\texorpdfstring{$\M_H$ 的拓扑}{M_H 的拓扑}}

%% ==== 原文 §3.6：PDF 第 62 页（印刷第 61 页）====
\label{sec:3.6}

在第~\ref{sec:3.5}~节中，我们考察了带抛物结构（parabolic structure）的丛的模空间（moduli space）$\M(\alpha;p)$ 的拓扑。这里我们将对希格斯丛（Higgs bundle）考虑类似的问题。

若简单地令 $\phi=0$，则希钦方程（Hitchin's equations）约化为平坦酉 $G$-丛（flat unitary $G$-bundle）的方程 $F=0$。根据上文提到的 Narasimhan 与 Seshadri 定理~\cite{ref51}，这类平坦丛的模空间同构于 $C$ 上稳定 $G$-丛的模空间，后者我们记作 $\M(G,C)$，并在上下文明确时简记为 $\M$。

\begin{yizhu}[物理图像补充]
令 $\varphi=0$ 得到的方程 $F=0$ 正是平坦酉 $G$-丛的方程，而把 $\M$ 嵌入 $\MH$ 相当于取 $\MH$ 的“零截面”。这一步的几何意义是：超 Kähler 结构的三个复结构各给出一种解释，在 $I$ 下 $\MH$ 是希格斯丛的模空间（本节的视角），在 $J$ 下它同构于平坦 $G_\C$-丛的模空间，这正是非阿贝尔霍奇对应（Donaldson--Corlette 与 Hitchin 的定理）；$\M$ 同时是两者的公共零点集。因此下文凡只涉及 $(\alpha,\beta,\gamma)$ 的拓扑陈述，在 $J$ 下都应读成关于平坦联络的陈述，而式~\eqref{eq:3.72} 的双有理等价也只是相对于复结构 $I$ 而言。
\end{yizhu}

在复结构（complex structure）$I$ 下审视，模空间 $\MH$ 参数化稳定对（stable pair）$(E,\varphi)$。我们可以通过遗忘 $\varphi$、只保留 $E$ 的全纯（holomorphic）类型，来定义 $\MH$ 的一个“叶状结构（foliation）”。对一般的稳定（或半稳定）对 $(E,\varphi)$，$E$ 是稳定（或半稳定）丛，而该叶状结构给出亚纯映射 $\psi:\M_H\to \M$（在 \cite{ref8} 中称为希钦第二纤维化（Hitchin's second fibration））。此映射的纤维是由 $\varphi\in H^0(C,K\otimes{\rm ad}(E))$ 参数化的线性空间；若 $E$ 稳定，该空间就是 $\M$ 在 $E$ 所确定的点处的余切空间（cotangent space）。此外，映射 $\psi$ 具有自然的截面（在复结构 $I$ 下全纯），因为我们可以把 $\M$ 嵌入 $\M_H$，作为 $\varphi=0$ 的希钦方程解构成的空间。于是在复结构 $I$ 下双有理地（birationally），$\MH$ 就是余切丛 $T^*\M$：
% 式 (3.72)
\setcounter{equation}{71}
\begin{equation}\label{eq:3.72}
  \MH\cong T^*\M.
\end{equation}

\begin{yizhu}[含糊点拨]
把遗忘 $\varphi$ 得到的结构叫作“叶状结构”只是隐喻：真正的全纯叶状结构要求纤维维数恒定、过每点有叶。此处的 $\psi$ 是亚纯映射，纤维 $H^0(C,K_C\otimes{\rm ad}(E))$ 的维数随 $E$ 变化——由 Riemann--Roch，$h^0(K_C\otimes{\rm ad}(E))=3g_C-3+h^1$，而 Serre 对偶给出 $h^1=h^0({\rm ad}(E))$，故 $E$ 稳定时纤维维数为 $3g_C-3$，$E$ 不稳定时更大；并且 $E$ 不稳定时 $\psi$ 根本没有定义。所以下文关于“希钦纤维化”的断言都应在稳定点附近理解，不能当作整体的叶状结构来用。
\end{yizhu}

更具体地，$\MH$ 以 $T^*\M$ 为稠密开集（dense open set）。与第~\ref{sec:3.5}~节所考虑的纤维化（fibration）有些类似，当 $C$ 的亏格（genus）$g_C$ 与 $G$ 的秩 $r$ 足够大时，这为 $\MH$ 的拓扑提供了一个有用的初步近似。例如，当 $g_C\geq 2$ 时，它足以用来讨论 $\MH$ 的二阶上同调（cohomology）。

\begin{yizhu}[含糊点拨]
这里的推理常被误读成“稠密开集蕴含 $H^2$ 同构”，其实不然：复余维 $\geq 1$ 的闭集虽不改变 $\pi_1$ 与 $H^1$，却可以改变 $H^2$（从复流形中挖去一点，$H^2$ 就会多出若干类），故单凭稠密性推不出 $H^2(\MH)\cong H^2(T^*\M)$。作者真正的依据是紧接着的论证：$\psi$ 的纤维是可缩的仿射空间，$\MH$ 作为仿射丛可收缩到 $\M$，于是 Leray 谱序列退化。两条理由结论相同，但只有后者是证明；“$g_C\geq 2$”与“秩 $r$ 足够大”这两个前提也正是为让不稳定点的集合余维足够高、从而不干扰低阶上同调。
\end{yizhu}

若参数 $\alpha$ 是 $\mathbb{L}$-正则的（$\mathbb{L}$-regular），则带分歧（ramification）的希格斯丛（ramified Higgs bundle）也有类似的结论。给定 $\M_{H,\mathbb{L}}(\alpha,\beta,\gamma;p)$ 中对应于稳定或半稳定的带分歧希格斯丛 $(E,\varphi)$ 的一个点，我们可以遗忘 $\varphi$，而把 $E$ 简单地看作带有由 $\alpha$ 确定的抛物结构（parabolic structure）的 $G_\C$-丛。一般地，该抛物丛（parabolic bundle）是稳定的，从而定义了 $\M_\mathbb{L}(\alpha;p)$ 中的一个点。于是我们得到亚纯映射 $\psi:\M_{H,\mathbb{L}}(\alpha,\beta,\gamma;p)\to \M_\mathbb{L}(\alpha;p)$。若 $\beta=\gamma=0$，情形与上面完全相同。映射 $\psi$ 的纤维是由希格斯场构成的空间，这些希格斯场具有幂零极点（nilpotent pole），取值于由 $(\mathbb{L},\alpha)$ 确定的抛物子群（parabolic subgroup）$\EUP$ 的幂幺根（unipotent radical）的李代数 $\mathfrak n$。这恰是 $\M_\mathbb{L}(\alpha;p)$ 的余切空间。此外，$\M_\mathbb{L}(\alpha;p)$ 可以全纯地嵌入 $\M_{H,\mathbb{L}}(\alpha,0,0;p)$，作为 $\varphi=0$ 的希钦方程解的空间。我们可以把这看作余切丛的零截面（zero section）。因此 $\M_{H,\mathbb{L}}(\alpha,0,0;p)$ 以 $\M_\mathbb{L}(\alpha;p)$ 的余切丛作为稠密开子空间包含之
% 式 (3.73)
\setcounter{equation}{72}
\begin{equation}\label{eq:3.73}
  \M_{H,\mathbb{L}}(\alpha,0,0;p)\cong T^*\M_\mathbb{L}(\alpha;p),
\end{equation}
且二者在复结构 $I$ 下双有理。

对 $\beta,\gamma\not=0$，这需要作一些修改。我们仍可遗忘 $\varphi$，从而定义亚纯映射
% 式 (3.74)
\setcounter{equation}{73}
\begin{equation}\label{eq:3.74}
  \psi:\M_{H,\mathbb{L}}(\alpha,\beta,\gamma;p)\to \M_\mathbb{L}(\alpha;p).
\end{equation}

然而，我们不能令 $\varphi$ 为零，因为其极点部分（polar part）的本征值（eigenvalue）由 $\beta$ 与 $\gamma$ 决定。因此映射 $\psi$ 没有全纯截面。于是 $\M_{H,\mathbb{L}}(\alpha,\beta,\gamma;p)$ 并不双有理于 $\M_\mathbb{L}(\alpha;p)$ 上的向量丛，而是双有理于一个“仿射丛（affine bundle）”。这意味着 $\psi$ 的纤维都是 $\C^N$ 的拷贝（对某个 $N$），但该纤维化的结构群（structure group）是仿射变换 $x\to ax+b$ 构成的群（而不只是线性变换 $x\to ax$）。与此相关，$\varphi$ 的极点部分不是幂零的，因此 $\varphi$ 并不代表 $\M_{H,\mathbb{L}}(\alpha;p)$ 的余切向量。但若 $\varphi$ 与 $\tilde\varphi$ 是两个希格斯场（具有相同的丛 $E$），则如式~\eqref{eq:3.26} 所示，它们的差 $\varphi-\tilde\varphi$ 的极点部分取值于 $\mathfrak n$，因而代表 $\MH(\alpha;p)$ 的余切向量。结论是：对一般的 $\beta$ 与 $\gamma$ 以及正则（regular）的 $\alpha$，$\MH(\alpha,\beta,\gamma;p)$ 含有一个稠密开集，它是 $\M(\alpha;p)$ 的余切丛的“仿射形变（affine deformation）”。

在作此陈述时，我们要求 $\alpha$ 是 $\mathbb{L}$-正则的，否则 $\M(\alpha;p)$ 会坍缩为维数更低的流形。然而，只要 $\beta$ 与 $\gamma$ 是一般的，无论 $\alpha$ 是否正则，$\MH(\alpha,\beta,\gamma;p)$ 的拓扑都与 $\alpha$ 无关~\cite{ref24}，因此对任意 $\alpha$ 与一般的 $\beta,\gamma$，我们得到了拓扑的一个粗略描述。当三元组 $(\alpha,\beta,\gamma)$ 非正则时，拓扑确实会在实余维（real codimension）三处发生改变，这一点我们马上说明。

\begin{yizhu}[易错点]
“拓扑与 $\alpha$ 无关”只能理解为微分同胚型不变，绝不是上同调类不变：式~\eqref{eq:3.86} 的 $[\omega_I/2\pi]=e\oplus(-\alpha^*)$ 显含 $\alpha$，并随 $\alpha$ 连续变化，不变的是整系数上同调群及其中的整格 $\Z\oplus\Lambda_{\rm wt}$。$\alpha$ 穿越非正则值时 $\MH$ 只作双有理变换（在超 Kähler 情形是 flop），故微分同胚型不变而辛结构改变。这与下一节的观察一致：$\alpha$ 只控制 Kähler 侧的消解参数 $r$，而 $\beta,\gamma$ 才控制复侧的形变参数 $\epsilon$（式~\eqref{eq:3.82}）。
\end{yizhu}

$\MH$ 与 $\M$ 之间关系的一个直接应用是：我们可以描述 $\MH$ 的二阶上同调，并确定描述辛形式（symplectic form）$\omega_I$ 的上同调类的公式~\eqref{eq:3.8}\ 中的未知常数。首先，$\MH$ 的二阶上同调同构于 $\M$ 的二阶上同调，因为作为仿射丛，$\MH$ 可收缩到 $\M$（除去余维（codimension）过高的部分——它们不影响二阶上同调，小亏格与小秩的特殊情形不在讨论之列）。因此
% 式 (3.75)
\setcounter{equation}{74}
\begin{equation}\label{eq:3.75}
  H^2(\MH;\Z)=\Z\oplus\Lambda_{wt}
\end{equation}
是 $G$ 的仿射权格（affine weight lattice），与 $H^2(\M;\Z)$ 的情形相同。此外，根据式~\eqref{eq:3.4}，$\omega_I=\psi^*(\omega)$，其中 $\omega$ 是 $\M$ 的通常辛形式，其上同调类已由式~\eqref{eq:3.51} 描述。于是我们可以直接借用式~\eqref{eq:3.51} 的结果，写出
% 式 (3.76)
\setcounter{equation}{75}
\begin{equation}\label{eq:3.76}
  \left[{\omega_I\over 2\pi}\right]=e\oplus (-\alpha^*).
\end{equation}

（此公式实际上对每个 $\mathbb{L}$ 都成立；它有意义，是因为 $\alpha^*$ 总取值于 $\Lambda_\wt\otimes_\Z\R$ 的 $\mathbb{L}$-不变部分。）由式~\eqref{eq:3.9}\ 与式~\eqref{eq:3.10} 可得
% 式 (3.77)
\setcounter{equation}{76}
\begin{equation}\label{eq:3.77}
  \left[{\omega_J\over 2\pi}\right]=0\oplus(-\beta^*),~~\left[{\omega_K\over 2\pi}\right]
=0\oplus(-\gamma^*).
\end{equation}
在讨论 $\MH(\alpha,\beta,\gamma;p)$ 的奇点之后，我们将能从这些公式得出一些有趣的结论。

\begin{yizhu}[含糊点拨]
作者用“$\MH$ 作为仿射丛可收缩到 $\M$”一语带过 $H^2(\MH)\cong H^2(\M)$，其实质是 Leray 谱序列：对 $\psi:\MH\to\M$ 有 $E_2^{p,q}=H^p(\M,R^q\psi_*\Z)\Rightarrow H^{p+q}(\MH;\Z)$，纤维是可缩的仿射空间，故 $R^q\psi_*\Z=0$（$q>0$），谱序列退化，$H^2(\MH)\cong H^2(\M)$ 由 $\psi^*$ 给出。严格而言 $\psi$ 只是亚纯映射（$E$ 不稳定处无定义），须把它限制在 $E$ 稳定的开子集上；该开子集的补集复余维不小于 $2$（作者所称“小亏格与小秩”以外的情形），故不影响 $H^2$。分解 $H^2=\Z\oplus\Lambda_{\rm wt}$ 中的 $\Z$ 由 $e$ 生成：$e$ 是行列式线丛 $\mathfrak L\to\M$ 的第一陈类的拉回（等价地，是 $\M$ 的辛形式 $\omega/2\pi$ 的上同调类的拉回），而 $-\alpha^*$ 中的负号与 $\alpha^*$ 的定义（经由二次型 $-\Tr$ 诱导的 $\mathfrak t\to\mathfrak t^\vee$）来自第~\ref{sec:3.5}~节的式~\eqref{eq:3.51}；与其它文献比较时须连同 $e$ 的归一化一起核对，不能只改符号。下一节中 $e\to e\oplus\mathfrak m^*$ 的“权平移”正是用来补偿 $\alpha^*\to\alpha^*+\mathfrak m^*$、使 $[\omega_I/2\pi]$ 保持不变（式~\eqref{eq:3.87}）；这是上同调层面一种 Tate 扭转式的平移，其整性依赖 $\mathfrak m^*\in\Lambda_{\rm wt}$。
\end{yizhu}

\bigskip\noindent{\it 特殊值 $(\alpha,\beta,\gamma)$ 处的奇点}

即使没有分歧，$G_\C$ 丛的模空间 $\M$ 与希格斯丛的模空间 $\M_H$ 也可能在对应于可约（reducible）丛或可约希格斯丛的点处有奇点（singularity）。这类奇点涉及 $C$ 上的整体行为，当 $C$ 的亏格大时出现在高余维处。

一点 $p\in C$ 处的抛物结构引入一类新的奇点，它只依赖于刻画分歧的参数。对丛而言这些参数是 $\alpha\in \TT$，对希格斯丛而言则是三元组 $(\alpha,\beta,\gamma)\in \TT\times {\mathfrak t}\times {\mathfrak t}$。这类奇点只依赖于 $p$ 处的局部行为，其余维与 $C$ 的亏格无关。在几何朗兰兹纲领（geometric Langlands program）的应用中，它们将比整体性的奇点重要得多。

对丛的情形，我们在第~\ref{sec:3.5}~节中描述了这类奇点，结果见式~\eqref{eq:3.67}\ 与式~\eqref{eq:3.71}。设 $\alpha$ 趋于一个值 $\bar\alpha$，它并非 $\mathbb{L}$-正则的，却是 $\mathbb{L}'$-正则的，其中 $\mathbb{L}'$ 是以 $\mathbb{L}$ 为真子群的列维群（Levi subgroup）。并设相应的抛物子群为 $\EUP$ 与 $\EUP'$。则当 $\alpha\to\bar\alpha$ 时，$\M_\mathbb{L}(\alpha;p)$ 的纤维是形如 $\mathbb{L}'/\mathbb{L}$（或 $\EUP'/\EUP$）的消没闭链（vanishing cycle）。纤维处处坍缩，故奇点的余维为零。图景如下
% 式 (3.78)
\setcounter{equation}{77}
\begin{equation}\label{eq:3.78}
  \begin{matrix}\EUP'/\EUP&\to & \M_\mathbb{L}(\alpha;p)\\
                           & &        \downarrow\\
                           & &      \M_{\mathbb{L}'}(\bar\alpha;p),\end{matrix}
\end{equation}
其中纤维当 $\alpha\to\bar\alpha$ 时坍缩。这是局部奇点，意即它只与 $\alpha$ 的取值以及 $p$ 点附近的行为有关，而与诸如寻找半稳定丛或可约平坦联络（flat connection）这样的整体问题无关。

那么这对带分歧的希格斯丛意味着什么？先看 $\beta$ 与 $\gamma$ 取最特殊的可能值、即 $0$ 时的情形。但我们从 $\alpha$ 的一个 $\mathbb{L}$-正则值出发。（关于这些问题的数学讨论见 \cite{ref61}，它从另一视角为下文许多陈述提供了依据。）如我们所讨论，$\M_{H,\mathbb{L}}(\alpha,0,0;p)$ 一般地是 $\M_\mathbb{L}(\alpha;p)$ 的余切丛。同样，若 $\mathbb{L}'$ 是真包含 $\mathbb{L}$ 的列维子群，且 $\bar\alpha$ 是 $\mathbb{L}'$-正则的，则 $\M_{H,\mathbb{L}'}(\bar\alpha,0,0;p)$ 一般地是 $\M_{\mathbb{L}'}(\bar\alpha;p)$ 的余切丛。从抛物丛过渡到带分歧希格斯丛，只需把式~\eqref{eq:3.78} 中的一切换成其对应的余切丛。于是当 $\alpha$ 接近 $\bar\alpha$ 时，$\M_{H,\mathbb{L}}(\alpha,0,0;p)$ 一般地纤维化于 $\M_{H,\mathbb{L}'}(\bar\alpha,0,0;p)$ 之上：
% 式 (3.79)
\setcounter{equation}{78}
\begin{equation}\label{eq:3.79}
  \begin{matrix}T^*(\EUP'/\EUP)=T^*(\mathbb{L}'/\mathbb{L})&\to& \M_{H,\mathbb{L}}(\alpha,0,0;p)\\
                                &   & \downarrow \\
                                & & \M_{H,\mathbb{L}'}(\bar\alpha,0,0;p). \end{matrix}
\end{equation}

我们用到了 $\EUP'/\EUP=\mathbb{L}'/\mathbb{L}$。当 $\alpha$ 趋于 $\bar\alpha$ 时，$\mathbb{L}'/\mathbb{L}$ 成为一个“消没闭链”，坍缩为一点。$\mathbb{L}'/\mathbb{L}$ 在其余切丛中的余维等于 $\mathbb{L}'/\mathbb{L}$ 的维数。因此这正是当 $\alpha$ 趋于非正则值 $\bar\alpha$ 时 $\M_{H,\mathbb{L}}(\alpha,0,0;p)$ 所获得的消没闭链的余维。在 $\alpha=\bar\alpha$ 处，消没闭链坍缩为一点，$\MH$ 获得一个奇点，该奇点的余维数加倍。我们称之为局部奇点，因为它只依赖于 $p$ 附近的局部行为（由参数 $(\alpha,\beta,\gamma)$ 决定），而不依赖于求解任何整体问题。例如，若 $\mathbb{L}=\TT$，且 $\bar\alpha$ 是一般的非正则值，则如第~\ref{sec:3.5}~节所见，$\mathbb{L}'/\mathbb{L}=\mathbb{CP}^1=S^2$。因此此情形下的消没闭链有实余维二。

由于 $\M_{H,\mathbb{L}}(\alpha,0,0;p)$ 是超 Kähler（hyper-Kähler）的，消没闭链附近的几何受到极强限制。几乎消没的闭链附近（$\alpha$ 接近 $\bar\alpha$ 时）高度弯曲的几何本身也必须是超 Kähler 的，$\M_{H,\mathbb{L}}(\alpha,0,0;p)$ 才可能是超 Kähler 的。可以显式地看清这是如何发生的。事实上，可利用纳姆方程（Nahm's equations）或希钦方程在 $T^*(\mathbb{L}'/\mathbb{L})$ 上构造一族超 Kähler 度规（metric）~\cite{ref34}。这是复李群 $\mathbb{L}'_\C$ 的余伴随轨道（coadjoint orbit）上超 Kähler 度规构造的一部分；我们将在第~\ref{sec:3.8}~节回顾这一构造。

与其动用这套完整机制，我们将较详细地考察 $\mathbb{L}=\TT$ 且 $\bar\alpha$ 为一般非正则值、从而 $\mathbb{L}'/\mathbb{L}=\mathbb{CP}^1$ 的情形。此时 $T^*(\mathbb{L}'/\mathbb{L})=T^*(\mathbb{CP}^1)$ 上相关的超 Kähler 度规族，是由 Calabi 以及 Eguchi 与 Hansen 最早构造的一族更熟悉的度规。

当 $\mathbb{CP}^1$ 的面积趋于零时，在 Calabi–Eguchi–Hansen 度规下 $T^*\mathbb{CP}^1$ 收敛到 $\R^4/\Z_2=\C^2/\Z_2$，即收敛到一个 $A_1$ 奇点。因此，在极限 $\alpha=\bar\alpha$ 处，$\M_{H,\mathbb{L}}(\alpha,0,0;p)$（由于已取 $\mathbb{L}=\TT$，它与 $\MH(\alpha,0,0;p)$ 相同）具有一族 $A_1$ 奇点。奇点集恰是嵌入在 $\M_H(\bar\alpha,0,0;p)$ 内部的超 Kähler 流形 $\M_{H,\mathbb{L}'}(\bar\alpha,0,0;p)$。这可由纤维化~\eqref{eq:3.79} 推得：当纤维 $T^*(\mathbb{L}'/\mathbb{L})$ 退化为具有孤立奇点的 $\C^2/\Z_2$ 时，$\M_H(\alpha,0,0;p)$ 的奇点集就成为该纤维化的一个截面，换言之，是 $\M_{H,\mathbb{L}'}(\bar\alpha,0,0;p)$ 的一个拷贝。下文我们还将给出得到这一结果的略有不同的途径（见式~\eqref{eq:3.82} 的讨论）。若 $\alpha$ 接近但不等于 $\bar\alpha$，或 $\alpha=\bar\alpha$ 但 $\beta$ 与 $\gamma$ 不恰好为零，则局部奇点被形变或消解，奇点附近的行为由 Calabi–Eguchi–Hansen 度规描述。这些论断可望在下文及第~\ref{sec:3.8}~节中变得清楚。

回顾关于 $\R^4/\Z_2$ 奇点的几个事实会有帮助。若选定一个复结构（在我们的应用中对应于 $\MH$ 的复结构 $I$），则 $\R^4/\Z_2$ 可描述为复奇点 $a_1^2+a_2^2+a_3^2=0$。它可以形变（deform）为光滑复流形
% 式 (3.80)
\setcounter{equation}{79}
\begin{equation}\label{eq:3.80}
  a_1^2+a_2^2+a_3^2=\epsilon.
\end{equation}
它也可以被消解而得到余切丛 $T^*\mathbb{CP}^1$，其中例外闭链（exceptional cycle）的面积为 $r$。此外，形变与消解可以同时进行。Calabi–Eguchi–Hansen 度规依赖于实参数 $r$ 以及复参数 $\epsilon$。$A_1$ 奇点恰在 $r=\epsilon=0$ 时出现；否则流形是光滑的。通过欧氏空间的超 Kähler 商（hyper-Kähler quotient）来构造 $\R^4/\Z_2$ 及其光滑形变~\cite{ref62}，可以非常自然地描述这一图景。或者，它是基于纳姆方程的构造的一个特例，我们将在第~\ref{sec:3.8}~节回顾该构造。

\begin{yizhu}[易错点]
把 $a_1^2+a_2^2+a_3^2=0$ 直接称作“$\R^4/\Z_2$”需要一点换算：$\Z_2$ 在 $\C^2$ 上的不变量环是 $\C[u,v]^{\Z_2}=\C[u^2,uv,v^2]$，其唯一关系为 $(u^2)(v^2)=(uv)^2$，即标准的 $A_1$ 形式 $UV=W^2$；而令 $U=a_1+ia_2$、$V=a_1-ia_2$、$W=ia_3$，便有 $UV=W^2$ 当且仅当 $a_1^2+a_2^2+a_3^2=0$。可见“$\R^4/\Z_2$ 的欧氏度规”与方程 $a_1^2+a_2^2+a_3^2=0$ 所带的代数复结构是两回事：前者固定渐近几何，后者只规定复结构；超 Kähler 结构的三个复结构 $I,J,K$ 才是把两者联系起来的数据。因此这里的 $a_i$ 不可当作 $\R^4$ 的线性坐标，它与 $r$ 分别对应形变与消解两个方向。
\end{yizhu}

现在把它应用于 $\MH$。在复结构 $I$ 中，$\alpha$ 是 Kähler 参数，如第~\ref{sec:2.2}~节的表格所述。因此变动 $\alpha$ 时，$\MH$ 只能通过双有理变换改变，且仅在出现消没闭链时才会改变（即三元组 $(\alpha,\beta,\gamma)$ 变为非正则时）。于是变动 $\alpha$ 给出 $A_1$ 奇点的爆破参数（blowup parameter）$r$。要看到复参数 $\epsilon$，则须变动 $\beta$ 与 $\gamma$。

\bigskip\noindent{\it $G=SU(2)$ 的分析}

为看清变动 $\beta$ 与 $\gamma$ 时会发生什么，为简单起见取 $G=SU(2)$。$\MH$ 中的点在复结构 $I$ 下对应于一个希格斯丛 $(E,\varphi)$。$\varphi$ 取值于 $V=H^0(C,K_C\otimes {\rm ad}(E)\otimes \CO(p))$；因此它是 $K_C\otimes {\rm ad}(E)$ 的一个截面，并可能在 $p$ 处有极点。空间 $V$ 的维数为 $3g_C$。$\varphi$ 受约束于
% 式 (3.81)
\setcounter{equation}{80}
\begin{equation}\label{eq:3.81}
  \Tr\,\varphi^2=\Tr\,\sigma^2 (dz/z)^2+\dots,
\end{equation}
其中 $\sigma=\half(\beta+i\gamma)$。

\begin{yizhu}[含糊点拨]
“空间 $V$ 的维数为 $3g_C$”是作者直接给出的，这里补出计算：$V=H^0(C,K_C\otimes{\rm ad}(E)\otimes\CO(p))$ 的秩为 $3$、次数为 $3(2g_C-2)+3=6g_C-3$（$\CO(p)$ 贡献次数 $3$，因为整个丛的秩是 $3$），Riemann--Roch 给出 $\chi=\deg+3(1-g_C)=3g_C$；而稳定 $E$ 上 $H^0({\rm ad}(E))=0$，由 Serre 对偶得 $h^1(V)=h^0({\rm ad}(E)(-p))=0$，故 $h^0=3g_C$。这既是系数 $a_i$ 恰有 $3g_C$ 个的理由，也是 $\dim\MH=6g_C-4=(3g_C-3)+3g_C-1$ 这一计数的依据；若 $E$ 不稳定则 $h^1>0$，维数还会更大。
\end{yizhu}

取 ${\rm ad}(E)$ 在 $p$ 处纤维的一组基 $b_i$，$i=1,2,3$，并归一化使 $\Tr\,b_ib_j=\delta_{ij}$。对 $i=1,2,3$，取元素 $\varphi_i\in V$，使其极点部分为 $\varphi_i\sim b_i(dz/z)$。再把 $\varphi_i$ 补充为 $V$ 的一组基 $\varphi_1,\dots,\varphi_{3g_C}$，使得 $i>3$ 的 $\varphi_i$ 在 $p$ 处没有极点。现在引入复参数 $a_i$，$i=1,\dots,3g$，并把 $\varphi$ 展开为 $\varphi=\sum_{i=1}^{3g_C}a_i\varphi_i$。$\MH(\alpha,\beta,\gamma;p)$ 由丛 $E$ 的选择（$3g_C-3$ 个参数）与系数 $a_i$（$3g_C$ 个参数）参数化，并受方程~\eqref{eq:3.81} 约束，故其维数为 $6g_C-4$。方程~\eqref{eq:3.81}\ 告诉我们
% 式 (3.82)
\setcounter{equation}{81}
\begin{equation}\label{eq:3.82}
  \sum_{i=1}^3a_i^2=\Tr\,\sigma^2.
\end{equation}
此方程描述 $A_1$ 奇点的形变，其中 $\Tr\,\sigma^2$ 扮演式~\eqref{eq:3.80} 中参数 $\epsilon$ 的角色。

\begin{yizhu}[原文疑误]
两处下标疑为原文笔误（已核对 PDF 图像，均照原文保留、未改公式）：(a) PDF 第 67 页（式~\eqref{eq:3.81} 之后）作“complex parameters $a_i$, $i=1,\dots,3g$”，但同句的和式上限为 $3g_C$，且空间 $V$ 的维数正是 $3g_C$，故 $3g$ 应为 $3g_C$（若把 $g$ 读作 $g_C$ 则无损）；(b) PDF 第 71--72 页（式~\eqref{eq:3.88} 之后）把余权族的指标范围写作 $i=1,\dots,n$、并把 $\Lambda^*$ 写成 $\oplus_{i=1}^n\Lambda_{{\rm cowt},i}$，而分歧点个数为 $s$，同段其余各处（$\sum_{i=1}^s u_i$、$\oplus_{i=1}^s\Lambda_{{\rm cort},i}$）均为 $s$，故这两处 $n$ 应为 $s$。另 (c) PDF 第 65 页“a singularity of whose codimension is twice as great”语法上应为“whose codimension”（或“of which the codimension”），是英文行文问题，译文按其本意处理。
\end{yizhu}

在 $\sigma=0$ 处，我们有望还原 $A_1$ 奇点，但此时必须记住控制 Kähler 结构的参数 $\alpha$。对一般的 $\alpha$，我们得到 $A_1$ 奇点的消解（resolution）。其理由很简单：当 $\varphi$ 的极点部分非零时，它（连同 $\alpha$ 的共轭类（conjugacy class）的指定）固定了 $\alpha$ 必须取的值。但当 $\varphi$ 没有极点时，$\alpha$ 的可能取值构成一族，且（由于 $\alpha$ 的共轭类固定）由 $\mathbb{CP}^1$ 参数化。由于 $\varphi$ 无极点的轨迹恰是 $A_1$ 奇点所在之处，把 $\alpha$ 的选择计入后，每个横截切片（transverse slice）中的 $A_1$ 奇点便被替换为一个 $\mathbb{CP}^1$ 的拷贝。这一操作就是奇点的消解。然而，若 $\alpha=0$，则 $\MH$ 确实会出现一族 $A_1$ 奇点。

\begin{yizhu}[易错点]
两个参数的角色不同，切勿混淆：$r$ 是消解参数（例外 $\mathbb{CP}^1$ 的面积，即 Kähler 参数），在复结构 $I$ 下由 $\alpha$ 一侧的数据控制；$\epsilon=\Tr\,\sigma^2$ 才是复形变参数，只由 $\beta,\gamma$ 控制。因此固定 $\beta,\gamma$ 而变动 $\alpha$ 只能作双有理变换（复结构 $I$ 不变，只改变辛/Kähler 结构），只有变动 $\beta,\gamma$ 才能看到 $A_1$ 奇点被形变。作者说“$\alpha$ 的共轭类固定时仍有由 $\mathbb{CP}^1$ 参数化的一族选择”，这里的 $\mathbb{CP}^1$ 实际上是 $p$ 处抛物结构的旗流形（对 $SU(2)$ 即 $G/B=\mathbb{CP}^1$），它只有在抛物子群真包含于 $G$ 时才非平凡；$\alpha=0$ 时 $P=G$，该自由度消失，故奇点不被消解——这与“一般 $\alpha$ 给出消解”并不矛盾。此外，式~\eqref{eq:3.82} 的归一化依赖 $\Tr\,b_ib_j=\delta_{ij}$（$b_i$ 为伴随丛在 $p$ 处纤维的基），$\Tr$ 取在基本表示上；换一种归一化会使右端出现因子 $2$。
\end{yizhu}

若令 $\alpha=\sigma=0$，则 $A_1$ 奇点的轨迹由 $a_1=a_2=a_3=0$ 给出。此条件保证 $\varphi$ 没有极点，从而 $(E,\varphi)$ 是不带分歧结构的普通希格斯丛。因此在本例中，$A_1$ 奇点的轨迹就是 $\M_H\subset \M_H(0,0,0;p)$。

这是一般地描述 $\M_{H,\mathbb{L}}(\alpha,\beta,\gamma;p)$ 的局部奇点的一个特例。若把 $(\alpha,\beta,\gamma)$ 取为三元组 $(\bar\alpha,\bar\beta,\bar\gamma)$，它非 $\mathbb{L}$-正则但为 $\mathbb{L}'$-正则，其中 $\mathbb{L}'$ 真包含 $\mathbb{L}$，则 $\M_{H,\mathbb{L}}(\bar\alpha,\bar\beta,\bar\gamma;p)$ 以 $\M_{H,\mathbb{L}'}(\bar\alpha,\bar\beta,\bar\gamma;p)$ 作为局部奇点的轨迹包含之。对给定的 $\mathbb{L}$，一切真包含 $\mathbb{L}$ 的可能 $\mathbb{L}'$ 都能出现，包括 $\mathbb{L}'=G$。

\subsection{仿射外尔群的作用}

%% ==== 原文 §3.7：PDF 第 68 页（印刷第 67 页）====
\label{sec:3.7}

为简单起见，本节开始时假设 $G$ 是单连通（simply-connected）的，且列维群为 $\mathbb{L}=\TT$。

参数 $(\alpha,\beta,\gamma)$ 一般取值于 $\TT\times \mathfrak t\times \mathfrak t$，并模去外尔群（Weyl group）${\mathcal W}$ 的作用。等价地，我们可以采取略有不同的观点，把 $(\alpha,\beta,\gamma)$ 看作取值于 $\mathfrak t\times \mathfrak t\times \mathfrak t$，并模去仿射外尔群（affine Weyl group）$\AffWeyl$ 的作用。我们回顾，$\AffWeyl$ 是 ${\mathcal W}$ 的一个扩张（extension），扩张核（kernel）为 $G$ 的余根格（coroot lattice）：
% 式 (3.83)
\setcounter{equation}{82}
\begin{equation}\label{eq:3.83}
  0\to \Lambda_{\rm cort}\to\AffWeyl\to {\mathcal W}\to 1.
\end{equation}

我们让 $\AffWeyl$ 作用在三元组 $(\alpha,\beta,\gamma)$ 上：对 $\alpha$ 按自然方式作用，而对 $\beta$ 与 $\gamma$ 则通过商群 ${\mathcal W}$ 作用。在此 $\AffWeyl$ 作用下，有
% 式 (3.84)
\setcounter{equation}{83}
\begin{equation}\label{eq:3.84}
  (\mathfrak t\times \mathfrak t\times \mathfrak t)/\AffWeyl=(\TT\times\mathfrak t\times \mathfrak t)/
{\mathcal W},
\end{equation}
这仅仅是因为对单连通的 $G$ 有 $\TT={\mathfrak t}/\Lambda_{\rm cort}$。对于被 $\mathcal W$ 作用的空间中的元素，若除单位元外 $\mathcal W$ 的任何元素都不使它保持不变，则称它是正则的；对被 $\AffWeyl$ 作用的空间中的元素同样如此。$(\alpha,\beta,\gamma)\in {\mathfrak t}^3$ 对 $\AffWeyl$ 的作用是正则的，当且仅当它到 $\TT\times \mathfrak t\times \mathfrak t$ 的投影对 $\Weyl$ 的作用是正则的。

现在，在模去 $\AffWeyl$ 之前，我们希望从 $\mathfrak t\times \mathfrak t\times \mathfrak t$ 中去掉那些使 $\MH(\alpha,\beta,\gamma;p)$ 出现奇点的点。我们在第~\ref{sec:3.6}~节分析了只依赖于 $p$ 附近行为的局部奇点。它们恰在三元组 $(\alpha,\beta,\gamma)\in \mathfrak t\times \mathfrak t\times \mathfrak t$ 对 $\AffWeyl$ 的作用非正则时出现（在第~\ref{sec:3.6}~节的讨论中，更自然的判据是等价的：到 $\TT\times\mathfrak t\times\mathfrak t$ 的投影按通常意义非正则）。这发生在实余维三处，因为被某个在 $\Weyl$ 中像非零的元素 $x\in \AffWeyl$ 固定，会分别对 $\alpha$、$\beta$、$\gamma$ 各施加一个非平凡条件。（而 $\Lambda_{\rm cort}$ 中的元素在 $\alpha$ 上自由作用，故根本没有不动点。）

\begin{yizhu}[易错点]
注意仿射外尔群的平移格与符号约定。对单连通 $G$ 有 $\TT=\mathfrak t/\Lambda_{\rm cort}$（余根格），式~\eqref{eq:3.83} 的扩张是半直积 $\AffWeyl=\Lambda_{\rm cort}\rtimes\Weyl$；若 $G$ 不单连通，则 $\TT=\mathfrak t/\Lambda_{\rm cochar}$（见本节末），而作用在单个连通分支上的仍是 $\Lambda_{\rm cort}$——把这两个格互换，是本节最常见的错误。平移元 $\mathfrak m\in\Lambda_{\rm cort}$ 的作用为 $\alpha\to\alpha+\mathfrak m$，而 $\beta,\gamma$ 只受商群 $\Weyl$ 的作用；正因平移不触及 $\beta,\gamma$，式~\eqref{eq:3.84} 才成立。反射 $s_i$ 的约定为 $s_i(\lambda)=\lambda-(\lambda,\alpha_i^\vee)\alpha_i$，并分别作用在 $\alpha,\beta,\gamma$ 上；若误把 $s_i$ 的平移部分也加到 $\beta,\gamma$ 上，就会得到错误的（余维更高的）奇点轨迹。余维计数正依赖于此：当 $x\in\AffWeyl$ 在 $\Weyl$ 中的像非平凡时，它同时固定 $\alpha$ 与 $\beta,\gamma$，各给出一个实条件，合计三个实条件，这正是“实余维三”的来源。
\end{yizhu}

此外还有整体奇点，它们出现在使希钦方程存在可约解的 $(\alpha,\beta,\gamma)$ 值处。即使我们已经适当调整三元组 $(\alpha,\beta,\gamma)$ 使这类解存在，只要 $C$ 的亏格大，它们在 $\MH$ 上仍只出现在高余维处。（这与局部奇点形成对照：后者一般是实余维四的 $A_1$ 奇点，如第~\ref{sec:3.6}~节所见。）这些奇点出现的 $(\alpha,\beta,\gamma)$ 值可以精确描述~\cite{ref24}。它们出现在 $\mathfrak t^3$ 中一族离散的、实余维三的仿射线性子空间（affine linear space）上。（实际上，分析表明：恰有一个抛物点时，对正则三元组 $(\alpha,\beta,\gamma)$ 不存在整体奇点；但有两个或更多抛物点时，可以出现整体奇点，而每个三元组仍为正则。）整体奇点处于实余维三的基本理由与局部奇点相同：$\M_H$ 的超 Kähler 本性。把 $\MH$ 视为其某个复结构下的复流形，要得到奇点，总须至少调节一个控制复结构的复参数与一个控制 Kähler 度规的实参数，合计三个实参数。

我们从 $\mathfrak t^3$ 中去掉所有使局部或整体奇点出现的余维三仿射线性子空间，把剩余部分记作 $\eurx$。由于 $\mathfrak t^3$ 是线性空间而去掉的部分余维为三，$\eurx$ 连通且单连通。当我们变动参数而不遇到奇点时，$\MH(\alpha,\beta,\gamma;p)$ 的拓扑不变。因此只要在 $\eurx$ 中变动参数，拓扑就不变。由于 $\eurx$ 连通，只要限制在 $(\alpha,\beta,\gamma)\in \eurx$ 内，流形 $\MH(\alpha,\beta,\gamma;p)$ 在拓扑上就与 $\alpha,\beta$ 以及 $\gamma$ 无关。这一陈述是 \cite{ref24} 的主要结果之一。

通过观察 $\AffWeyl$ 在 $\eurx$ 上自由作用（因为所有三元组 $(\alpha,\beta,\gamma)\in \eurx$ 都是正则的），我们可以知道更多。于是 $\eurx/\AffWeyl$ 是光滑流形，其基本群为 $\AffWeyl$。$\MH(\alpha,\beta, \gamma;p)$ 的上同调作为 $\eurx/\AffWeyl$ 上某个平坦丛（flat bundle）的纤维而变化。取该平坦丛的单值性（monodromy），就得到 $\AffWeyl$ 在 $\MH(\alpha,\beta,\gamma;p)$ 的上同调上的作用。类似地，$\AffWeyl$ 作用于此空间的例如 $K$-理论（$K$-theory）上。

\begin{yizhu}[易错点]
务必区分两件事：$\AffWeyl$ 作用在参数空间 $\eurx$ 上（这是构造），而它在 $H^*(\MH)$ 上的“作用”是由平坦丛的单值性得到的表示 $\rho:\AffWeyl\to{\rm Aut}\,H^*(\MH;\Z)$，并非流形上的群作用——$\MH$ 上并没有 $\AffWeyl$ 作用，二者的关系只是“$\MH$ 的参数族以 $\eurx/\AffWeyl$ 为底”。这一区分有可检验的后果：$\rho$ 只在共轭意义下确定，换基点或改走另一条道路会得到共轭的表示，而真正的群作用不会有这种含糊。$K$-理论的情形同理，得到的是取值于 $K^0$ 自同构群中的单值性表示，不能与上同调上的表示混为一谈。
\end{yizhu}

\begin{yizhu}[含糊点拨]
“$\eurx$ 连通且单连通”“拓扑不随参数改变”两句话省略了技术前提。第一，$\mathfrak t^3$ 去掉若干余维三的仿射线性子空间后仍连通且单连通：实余维三的闭集不改变 $\R^n$（$n\geq 3$）的连通性与单连通性，可用横截球面 $S^2$ 的连通性看出。第二，参数族的局部平凡性用的是 Ehresmann 引理，需要对一族紧致或固有（proper）的光滑流形——此处应经由希钦纤维化或超 Kähler 商的固有性来保证；仅凭“奇点不出现”并不自动给出微分同胚。第三，$\AffWeyl$ 在 $\eurx$ 上自由，等价于每个三元组正则（稳定子平凡），因此 $\eurx/\AffWeyl$ 是光滑流形且基本群为 $\AffWeyl$；若 $\eurx$ 只是道路连通而非单连通，单值性给出的将是 $\AffWeyl$ 的某个扩张的作用，作者此处用到了单连通性来排除这种可能。
\end{yizhu}

这一结果多少类似于外尔群的施普林格表示（Springer representation）~\cite{ref26}；如 \cite{ref31} 所建议，通过变动余伴随轨道及其 Slodowy 切片上超 Kähler 度规的参数，也可以自然地理解这些表示~\cite{ref35}。（我们在第~\ref{sec:3.8}~节回顾这一框架。）施普林格表示已被推广为仿射赫克代数（affine Hecke algebra）在等变 $K$-理论（equivariant $K$-theory）上的作用~\cite{ref27,ref28}，最近又被推广为仿射辫群（affine braid group）在施普林格消解（Springer resolution）上的层（sheaf）的某些导出范畴（derived category）上的作用~\cite{ref29}。关于其中部分结果，可参看 Chriss 与 Ginzburg 的书~\cite{ref30}。$\MH$ 的上同调或 $K$-理论上的仿射外尔群作用，将进一步丰富为仿射辫群在 $A$-膜（A-brane）或 $B$-膜（B-brane）范畴上的作用，我们将在第~\ref{sec:4.5}~节看到这一点。至于能否在 $\MH$ 的语境中看到仿射赫克代数，若要尝试，就令 $\beta=\gamma=0$，此时 $\MH$ 容许 $\C^*$ 的作用，从而可以定义其等变上同调或 $K$-理论。它们很可能容许仿射赫克代数的作用，这将改进涉及 $\MH$ 的“整体”问题与涉及复余伴随轨道的“局部”问题之间的类比。

\begin{yizhu}[易错点]
这一段叠着三个层次，不可混。其一是仿射外尔群 $\AffWeyl$ 本身：它由参数空间 $\eurx$ 上的单值性得到，作用在 $H^*$ 与 $K$-理论上，这是本节的定理。其二是仿射赫克代数 $\mathcal H_{\rm aff}$：它是 $\C[\AffWeyl]$ 的 $q$-形变（$q\to1$ 时退化），Kazhdan 与 Lusztig 证明它在施普林格情形的等变 $K$-理论上作用，其标准基（KL 基）的乘法常数就是 Kazhdan--Lusztig 多项式。其三是仿射辫群：它是 $\AffWeyl$ 的辫子化，作用在导出范畴（$A$/$B$-膜范畴）上，属范畴化层面。作者对 $\MH$ 上是否存在 $\mathcal H_{\rm aff}$ 作用只说“很可能”（may well admit），那是猜想，不可当作已证结论。
\end{yizhu}

\bigskip\noindent{\it 例}

仿射外尔群作用在 $\MH$ 各个维数的上同调上。不过，若限制到二维上同调（我们在式~\eqref{eq:3.75} 中描述过它），则可以显式地描述这一作用：
% 式 (3.85)
\setcounter{equation}{84}
\begin{equation}\label{eq:3.85}
  H^2(\MH(\alpha,\beta,\gamma;p);\Z)=\Z\oplus \Lambda_{\rm wt}.
\end{equation}

右端是 $G$ 的仿射权格，因此自然地接受 $\AffWeyl$ 的作用。为了给这个显然的猜想——$\AffWeyl$ 确实以这种自然方式作用在 $H^2(\MH)$ 上——提供依据，我们利用关于辛形式 $\omega_I$ 的上同调类的结果~\eqref{eq:3.76}：
% 式 (3.86)
\setcounter{equation}{85}
\begin{equation}\label{eq:3.86}
  \left[{\omega_I\over 2\pi}\right]=e\oplus (-\alpha^*).
\end{equation}

\begin{yizhu}[含糊点拨]
“仿射权格”指 $\Z\oplus\Lambda_{\rm wt}$：多出的 $\Z$ 对应仿射李代数的 level（此处 level 固定为 $1$），$\Lambda_{\rm wt}$ 是 $G$ 的权格。$\Weyl\subset\AffWeyl$ 平凡地作用在 $\Z$ 上（保 level），这正是下句“$\Weyl$ 平凡地作用在 $\Z$ 上”的内容；而平移元 $\mathfrak m\in\Lambda_{\rm cort}$ 经附录 A 的映射 $u\mapsto u^*$ 给出 $\mathfrak m^*\in\Lambda_{\rm wt}$，作用在第二个直和项上。因此式~\eqref{eq:3.87} 与式~\eqref{eq:3.85} 的相容性是显然的：$e$ 与 $\mathfrak m^*$ 分处两个直和项、互不干扰；但若把 $\mathfrak m^*$ 错放进 $\Z$ 这一项，$e$ 的整性就会被破坏。
\end{yizhu}

$\omega_I$ 的上同调类必须在 $\AffWeyl$ 对 $\alpha$ 与对 $H^2(\MH(\alpha,\beta,\gamma;p);\Z)$ 的联合作用下保持不变。这唯一地确定了 $\AffWeyl$ 在 $H^2(\MH(\alpha,\beta,\gamma;p);\Z)$ 上的作用就是其在仿射权格上的自然作用。这意味着子群 $\Weyl\subset \AffWeyl$ 平凡地作用在 $\Z$ 上，并按通常方式作用在 $\Lambda_{\rm wt}$ 上。而 $\mathfrak{m}\in \Lambda_{\rm cort}$ 的作用为
% 式 (3.87)
\setcounter{equation}{86}
\begin{equation}\label{eq:3.87}
  e\to e\oplus {\mathfrak{m}}^*,
\end{equation}
而在 $\Lambda_{\rm wt}$ 上平凡地作用。显然，\eqref{eq:3.86}\ 的右端在此变换连同 $\alpha\to\alpha+\mathfrak{m}$ 之下保持不变。

\bigskip\noindent{\it 多个分歧点}

我们可以容易地把这一点推广到有多个分歧点 $p_1,\dots,p_s$ 的情形。每个这样的点对应一个三元组 $(\alpha_i,\beta_i,\gamma_i)\in {\mathfrak t}^3$，并带有各自的 $\AffWeyl$ 作用。于是群 $(\AffWeyl)^s$ 作用在这 $s$ 个三元组的集合上，后者取值于 $(\mathfrak t^3)^s={\mathfrak t}^{3s}$。

当诸三元组 $(\alpha_i,\beta_i,\gamma_i)$ 取值于适当的参数空间 $\eurx_s$ 时，相应的带分歧希格斯丛的模空间 $\MH(\alpha_1,\beta_1,\gamma_1,p_1; \dots;\alpha_s,\beta_s,\gamma_s,p_s)$ 是光滑流形。该空间由 $\mathfrak t^{3s}$ 去掉某些余维三的仿射线性子空间得到。为避免局部奇点，必须要求每个三元组 $(\alpha_i,\beta_i,\gamma_i)$ 分别正则。为避免整体奇点，还须去掉 \cite{ref24} 中描述的某些额外的仿射线性子空间。

群 $(\AffWeyl)^s$ 在 $\eurx_s$ 上自由作用，因为在定义 $\eurx_s$ 时我们要求每个三元组分别正则。因此，按与前面相同的逻辑，我们得到 $(\AffWeyl)^s$ 在 $\MH$ 的上同调、$K$-理论等上的作用。

与前面一样，若专门考察二维上同调，则可以显式地描述这一作用，该上同调已在 \eqref{eq:3.48} 中描述：
% 式 (3.88)
\setcounter{equation}{87}
\begin{equation}\label{eq:3.88}
  H^2(\M(\alpha_1,p_1;\dots;\alpha_s,p_s);\Z)=\Z\oplus\left(\oplus_{i=1}^s
\Lambda_{wt,i}\right).
\end{equation}

通常外尔群的 $s$ 个拷贝作用在 $\Lambda_{wt}$ 的 $s$ 个拷贝上。而格（lattice）的作用是 \eqref{eq:3.87} 的推广：
% 式 (3.89)
\setcounter{equation}{88}
\begin{equation}\label{eq:3.89}
  e\to e\oplus\left(\oplus_{i=1}^s {\mathfrak{m}}_i^*\right).
\end{equation}

这确实是 $(\AffWeyl)^s$ 的正确作用，这一点实际上可由 \eqref{eq:3.86} 的类比推出，即 $[\,\omega_I/2\pi]=e\oplus\left(\oplus_i(-\alpha^*_i)\right)$。

人们可能会对这些结果感到困惑，因为一个稍大的群也能作用在格~\eqref{eq:3.88} 上。与其用一族余根 $\mathfrak{m}_i$ 把 $e$ 平移 $\oplus_{i=1}^s {\mathfrak{m}}^*_i$，为什么不直接任取一组权 $w_1,\dots,w_s\in \Lambda_{\rm wt}$ 而把 $e$ 平移 $\oplus_{i=1}^s w_i$，不必操心 $w_i$ 是否形如某个余根 $\mathfrak{m}_i$ 的 ${\mathfrak{m}}^*_i$？事实上，在多个分歧点的情形，一个比这更小、但比我们迄今所描述的更大的群，确实自然地作用在 $\MH( \alpha_1,\dots,p_s)$ 的上同调上。根据 \eqref{eq:3.63}\ 与 \eqref{eq:3.64}，对任意一族余权（coweight）$u_i$，$i=1,\dots,n$，只要 $\sum_{i=1}^su_i$ 是余根，就存在对称性 $\alpha_i\to \alpha_i+u_i$。（我们是对抛物丛证明该对称性的，但同样的推理适用于带分歧的希格斯丛。）因此 $\MH(\alpha_1,\dots,p_s)$ 的上同调必须容许群 $\Lambda^*\rtimes (\Weyl)^s$ 的作用，其中 $\Lambda^*$ 是 $\oplus_{i=1}^n \Lambda_{{\rm cowt},i}$ 中由元素 $u_1\oplus\dots \oplus u_s$（所有 $u_i\in \Lambda_{\rm cowt}$ 且 $\sum_iu_i\in \Lambda_{\rm cort}$）构成的子格。由于 $\Lambda^*$ 以 $\oplus_{i=1}^s\Lambda_{{\rm cort},i}$ 为有限指数子格，而前者的作用我们已描述过，故其作用必由同一公式给出：
% 式 (3.90)
\setcounter{equation}{89}
\begin{equation}\label{eq:3.90}
  e\to
e\oplus\left(\oplus_{i=1}^s u^*_i\right).
\end{equation}
这是合理的，因为映射 $u\to u^*$ 确实把 $\Lambda_{\rm cowt}$ 映入 $\Lambda_{\rm wt}$，如附录 A 所述。

\begin{yizhu}[易错点]
为什么不能“任取一组权”？因为单点平移 $\alpha_i\to\alpha_i+u_i$ 会按式~\eqref{eq:2.6} 把该点的特征类 $\xi$ 平移 $y_i[D_i]$；只有当 $\sum_i u_i$ 落在余根格 $\Lambda_{\rm cort}$ 中时，这些平移的整体效应才与整体规范变换等价，模空间才留在同一拓扑型上——这正是式~\eqref{eq:3.63} 与~\eqref{eq:3.64} 的条件。所以 $\Lambda^*$ 是 $\oplus_i\Lambda_{{\rm cowt},i}$ 中被“总平移属于 $\Lambda_{\rm cort}$”截出的子格：它含 $\oplus_i\Lambda_{{\rm cort},i}$ 而含于 $\oplus_i\Lambda_{{\rm cowt},i}$，两者之商即 $\Lambda_{\rm cowt}/\Lambda_{\rm cort}$，故前者在后者中指数有限。
\end{yizhu}

\bigskip\noindent{\it 非单连通的 $G$}

若 $G$ 不是单连通的，会怎样？此时 $\alpha$ 取值于 $\TT=\mathfrak t/ \Lambda_{\rm cochar}$，看起来作用在 $\MH$ 的上同调或 $K$-理论上的群应该是 $\Lambda_{\rm cochar}\rtimes \Weyl$。

这是否正确，取决于所说的究竟是什么。当 $G$ 不单连通时，$\MH$ 有 $\# \pi_1(G)$ 个连通分支，由特征类（characteristic class）$\xi(E)$ 的值标记，后者度量把 $E$ 提升为结构群为单连通群 $\bar G$ 的丛的障碍。$\Lambda_{\rm cochar}$ 的作用置换这些分支，如式~\eqref{eq:2.6} 所讨论。作用在给定分支的上同调上的子群是 $\Lambda_{\cort}$。因此，若只关注一个固定分支，则在单个抛物点的情形，作用在上同调上的群为 $\AffWeyl=\Lambda_{\cort}\rtimes \Weyl$，与 $G$ 单连通时完全一样。类似地，对 $s$ 个抛物点，无论 $G$ 是否单连通，作用在 $\MH$ 的单个分支的上同调上的群都是 $\Lambda^*\rtimes (\Weyl)^s$。但若要计入置换各分支的变换，则需用更大的群。

\begin{yizhu}[含糊点拨]
分支数恰为 $\#\pi_1(G)$ 的理由是：特征类 $\xi$ 取值于 $H^2(C;\pi_1(G))$，而 $C$ 是实二维紧黎曼面，由万有系数定理 $H^2(C;A)\cong{\rm Hom}(H_2(C),A)\oplus{\rm Ext}(H_1(C),A)\cong A$（因 $H_2(C)=\Z$、$H_1(C)=\Z^{2g_C}$ 自由），故 $\xi$ 只取 $\#\pi_1(G)$ 个值；对 $G=SO(3)$ 这就是 $w_2(E)\in H^2(C;\Z_2)=\Z_2$ 的两种取值。$\Lambda_{\rm cochar}$ 按式~\eqref{eq:2.6} 置换这些分支，其稳定子恰为 $\Lambda_{\rm cort}$，于是单个分支上作用的仍是 $\Lambda_{\rm cort}\rtimes\Weyl=\AffWeyl$——这正是“与 $G$ 单连通时完全一样”的含义。
\end{yizhu}

\bigskip\noindent{\it $\mathbb{L}$ 型分歧}

我们可以用类似方式考虑一个带有由任意列维子群 $\mathbb{L}$ 标记的奇点的点 $p$。标记该点的参数 $(\alpha,\beta,\gamma)$ 在 $\mathbb{L}$ 下不变，而局部奇点恰在此三元组为 $\mathbb{L}$-正则时才被避免。自然作用的群是与 $\mathbb{L}$ 对易的 $\AffWeyl$ 的子群。我们称此群为 ${\mathcal W}_{{\rm aff},\mathbb{L}}$。

与上面相同的推理表明，$\Weyl_{{\rm aff},\mathbb{L}}$ 作用在 $\M_{H,\mathbb{L}}$ 的上同调（或 $K$-理论等）上。在 $\mathbb{L}=G$ 的极端情形，此陈述变得平凡，因为此时 $\Weyl_{{\rm aff},\mathbb{L}}$ 是平凡群。

\begin{yizhu}[易错点]
“与 $\mathbb{L}$ 对易的子群”不可按字面读成“与 $\mathbb{L}$ 逐元素对易的 $\AffWeyl$ 元素之集”：$\mathbb{L}$ 并不是 $\AffWeyl$ 的子群，若照字面算出 $\Lambda_{\rm cort}(\mathbb{L})\rtimes\Weyl_{\mathbb{L}}$，在 $\mathbb{L}=G$ 时将得到整个 $\AffWeyl$，与作者“此时该群平凡”的结论相反。作者的实际用法应读作“保持该型奇异数据的单值性群”：$\mathbb{L}=G$ 时 $p$ 处没有分歧数据，$\M_{H,\mathbb{L}}$ 就是普通的 $\MH$，没有可环绕的奇异轨迹，故该群平凡；这也与上文“局部奇点仅在非 $\mathbb{L}$-正则时出现”的判据一致。
\end{yizhu}

\bigskip\noindent{\it 推广到四维规范场论与西格玛模型}

在几何朗兰兹纲领的应用中，我们真正关心的并非流形 $\MH$，而是以 $\MH$ 为目标空间的二维西格玛模型（sigma model）。该西格玛模型作为四维规范场论（gauge theory）的低能近似出现。比经典流形是否光滑更重要的，是西格玛模型是否光滑。当西格玛模型的谱（spectrum）、关联函数（correlation function）等随参数光滑变化时，我们称它是光滑的。

当 $\MH$ 光滑时，西格玛模型当然光滑。然而，即使经典上 $\MH$ 有奇点，西格玛模型也可能仍然光滑。通常，在具有 $(4,4)$ 超对称（supersymmetry）的二维西格玛模型中，要得到量子理论（quantum theory）的奇点，必须调节四个参数，而不只是三个。三个参数控制经典几何，第四个参数控制西格玛模型的一个 $\theta$ 角（theta angle）。

如我们在第~\ref{sec:2.3}~节所学，除参数 $(\alpha,\beta,\gamma) \in \TT\times \mathfrak t\times \mathfrak t$ 外，规范场论还依赖于另一个参数，即 $\theta$ 角 $\eta\in \,^L\neg\TT$。因此标记带分歧结构的点的参数实际上是一个四元组\footnote{若愿意，我们可以把 $\alpha$ 与 $\eta$ 提升为取值于 $\mathfrak t$ 的量。此时所需模去的群不是仿射外尔群，而是 $\Weyl$ 被一对格的乘积所作的一个扩张，这对格通过平移 $\alpha$ 与 $\eta$ 起作用，如我们在第~\ref{sec:4.5}~节所讨论。} $(\alpha,\beta,\gamma,\eta)\in \TT\times \mathfrak t\times \mathfrak t\times ^L\neg\TT$，并在外尔群对全部四个变量的作用下等价。这是合理的，因为 $G$ 与 $^LG$（对偶群，dual group）有相同的外尔群！

所谓以 $\MH$ 为目标空间的西格玛模型的局部奇点，是指其位置只依赖于标记单个分歧点的参数 $(\alpha,\beta,\gamma,\eta)$ 的奇点。整体奇点则是指其位置依赖于两个或更多点的参数的奇点。

我们回顾：被 $\Weyl$ 作用的空间中的点，若没有 $\Weyl$ 的非平凡元素使其固定，则称为正则的。为使讨论简单，我们从单个分歧点的情形开始，并利用此时 $\MH$ 只在三元组 $(\alpha,\beta,\gamma)\in \TT\times\mathfrak t\times \mathfrak t$ 非正则时才有奇点这一事实。由此可知，若 $(\alpha,\beta,\gamma)$ 正则，则西格玛模型光滑。然而，规范场论还有 $S$-对偶（$S$-duality）变换 $\tau\to -1/\tau$，它交换 $\alpha$ 与 $\eta$，如我们在第~\ref{sec:2.4}~节所讨论。西格玛模型的光滑性必须在此变换下不变，因此得知：若三元组 $(\eta,\beta,\gamma)$ 正则，则西格玛模型光滑。

更一般地，我们有一个无限的离散对偶群 $\Gamma$，按 \eqref{eq:2.53} 作用在 $\alpha,\eta$ 上。若 $G$ 是单边（simply-laced）的，或我们限制到 $\Gamma$ 的指数为二的子群，则 $\beta,\gamma$ 按乘以正实数的方式变换（注意 \eqref{eq:2.24} 之后的讨论）。这样的重标度不影响三元组 $(\alpha,\beta,\gamma)$ 是否正则这一问题。例如对 $G=E_8$，对偶群是 $SL(2,\Z)$，它含有把 $\alpha$ 映为 $m\alpha+n\eta$ 的变换，其中 $m,n$ 为任意互素整数。因此若 $(m\alpha+n\eta,\beta,\gamma)$ 正则，则西格玛模型光滑。对所有 $m,n$ 均如此这一事实意味着：实际上，若四元组 $(\alpha,\beta,\gamma,\eta)$ 正则，则西格玛模型光滑。若把 $SL(2,\Z)$ 换成同余子群（congruence subgroup），仍会得到同样结论。对任意 $G$，对偶群都含有一个按上述方式作用的 $SL(2,\Z)$ 的同余子群，因此总有：若四元组 $(\alpha,\beta,\gamma,\eta)$ 正则，则西格玛模型光滑。

\begin{yizhu}[物理图像补充]
物理上这四个参数的地位并不相同。$\alpha,\beta,\gamma$ 是四维规范场论中刻画带分歧的面算符（surface operator）的数据，在超 Kähler 的三个复结构 $I,J,K$ 下分别与三个辛形式 $\omega_I,\omega_J,\omega_K$ 的 Kähler 类相对应（因此在复结构 $I$ 中只有 $\alpha$ 是 Kähler 参数）；而 $\eta$ 是四维理论的 $\theta$ 角，属于量子参数（二维西格玛模型的 B-场），在经典模空间中不可见。这正解释了“经典奇点需调节三个参数、量子奇点需调节四个参数”：$(4,4)$ 超对称的二维西格玛模型要出现量子奇点，必须同时调节三个经典几何参数与一个 $\theta$ 角。$S$-对偶 $\tau\to-1/\tau$ 交换 $\alpha$ 与 $\eta$，更一般地，对偶群 $\Gamma$ 按 \eqref{eq:2.53} 把二者混合，因此光滑性判据必须在四元组的任意 $\Gamma$-像上同时成立；再借助 $\Gamma$ 含有 $SL(2,\Z)$ 的同余子群，便得到不依赖于对偶的“四元组正则”判据。从物理上看，$\theta$ 角的平移与电磁荷（威尔逊线／'t Hooft 线）的平移相联系（Witten 效应），这正是本节中仿射外尔群平移作用的一个来源。
\end{yizhu}

对于在多个点 $p_i$ 处分歧的情形，我们同样可以考虑整体奇点；可以简单地把它定义为位置依赖于与多于一个点相关的参数 $(\alpha_i,\beta_i,\gamma_i,\eta_i)$ 的奇点。现在可以利用对偶对称性证明：\cite{ref24} 中找到的条件可以再加上一个涉及诸 $\eta$ 的附加条件，于是整体奇点现在出现在实余维四处。
